\documentclass[11pt]{article}
\usepackage[T1]{fontenc}
\usepackage{lmodern}
\usepackage{amsmath,amssymb,amsthm}
\usepackage[margin=1in]{geometry}
\usepackage{microtype}
\usepackage{xurl}
\usepackage[hidelinks,breaklinks]{hyperref}
\hypersetup{pdftitle={A certified coordinate-gauge obstruction for causal-order reconstruction in a 2+1 diamond},pdfauthor={Daniel Eric Fredriksen},pdfkeywords={causal order, conformal diamond, coordinate gauge, Lean, isometry}}
\newtheorem{theorem}{Theorem}[section]
\newtheorem{proposition}[theorem]{Proposition}
\newtheorem{lemma}[theorem]{Lemma}
\newtheorem{corollary}[theorem]{Corollary}
\theoremstyle{remark}
\newtheorem{remark}[theorem]{Remark}
\newcommand{\R}{\mathbb R}
\newcommand{\D}{\mathcal D}
\newcommand{\K}{\mathcal K_3}
\newcommand{\TV}{\operatorname{TV}}
\newcommand{\Otwo}{\mathrm O(2)}
\newcommand{\dd}{\mathrm d}
\newcommand{\code}[1]{{\small\ttfamily\def\_{\textunderscore\allowbreak}#1}}
\title{A certified coordinate-gauge obstruction for causal-order reconstruction in a $2+1$ diamond}
\author{Daniel Eric Fredriksen\\Quantyra Inc.\\\url{https://github.com/Quantyra/quantyra-null-cone-lean}}
\date{8 October 2026\\Version 0.1.0 --- unpublished manuscript}

\begin{document}
\maketitle
\begin{abstract}
Finite chronological-order samples are invariant under a smooth change of
coordinates that transports both chronology and sampling measure. We give an
explicit Lean-certified example showing why quotienting density coefficients
only by spatial rotations and reflections does not remove this freedom.
On the unit $2+1$ Minkowski diamond, the standard conformal diamond flow
transports normalized flat volume to a nonconstant smooth density between
$1/2$ and $3/2$ with Euclidean Lipschitz constant at most two.
Every finite labeled and unlabeled directed-order law agrees with the flat
law, while the coordinate supremum distance remains positive after every
spatial orthogonal transformation. Nevertheless, the volume-normalized
Lorentzian metrics are smoothly isometric, with an explicit pullback identity.
We prove each step, identify the flow with its established conformal-field-theory
formula, and give a statement-to-source map for the accepted formalization.
The contribution is an explicit certified instance and an audit of the
proposed coordinate gauge. The flow and the general invariance principle
are established; no new physical nonidentifiability or priority claim is made.
\end{abstract}

\noindent\textbf{Status:} separate unpublished technical manuscript; no DOI assigned.\\
\textbf{Artifacts:} source, Lean proofs and GCP acceptance evidence in the repository above.\\
\textbf{License:} manuscript CC-BY-4.0; proof software Apache-2.0.

\section{Question, result and scope}
An inverse problem based on finite causal orders must specify which
coordinate descriptions represent the same object. In the reconstruction
paper~\cite{Fredriksen}, a restricted $1+1$ density class fixes uniform
null-coordinate marginals. A tentative higher-dimensional extension would
instead compare density functions on a fixed diamond modulo spatial
rotations and reflections. The present example shows that this smaller
symmetry group leaves an observable ambiguity inside a smooth, uniformly
bounded class.

The obstruction uses an established conformal flow. Its mathematical value
here is the complete passage from a Cartesian formula to an admissible
density, actual probability laws, a positive coordinate discrepancy and an
explicit isometry, together with machine-checked evidence. Section~\ref{sec:prior}
compares the construction with primary sources. We do not treat certification
as a proof of originality.

\subsection{The fixed model}
Write $p=(t,z)=(t,x,y)$, $r=|z|_2$, and set
\[
 \D=\{(t,z): |t|+|z|_2<1\},\qquad
 \eta=-\dd t^2+\dd x^2+\dd y^2.
\]
An overline denotes closure. The underlying coordinate norm is the
Euclidean norm on $\R^3$. The chronological relation is the genuine
Lorentzian relation
\begin{equation}\label{eq:chronology}
 p\ll q\quad\Longleftrightarrow\quad q_t-p_t>|q_z-p_z|_2.
\end{equation}
The diamond is convex. Thus the straight segment realizes every relation
in~\eqref{eq:chronology}; conversely, integrating the future timelike cone
inequality along a timelike curve gives the same endpoint condition.
Consequently this is also the intrinsic chronology of the diamond.

Its flat volume and normalized reference measure are
\begin{equation}\label{eq:volume}
 V=\int_{-1}^{1}\pi(1-|t|)^2\,\dd t=\frac{2\pi}{3},
 \qquad \mu_0=V^{-1}\,\dd t\,\dd x\,\dd y\big|_{\D}.
\end{equation}
Let $\K$ consist of real functions smooth on an open neighborhood of
$\overline\D$ such that
\begin{equation}\label{eq:class}
 \frac12\leq\rho\leq\frac32,\qquad
 |\rho(p)-\rho(q)|\leq2\|p-q\|_2\quad(p,q\in\overline\D),
 \qquad \int\rho\,\dd\mu_0=1.
\end{equation}
For $\rho\in\K$, put $\mu_\rho=\rho\mu_0$ and
\begin{equation}\label{eq:metric}
 g_\rho=V^{-2/3}\rho^{2/3}\eta.
\end{equation}
The Lorentzian volume density of this three-dimensional metric is
$V^{-1}\rho$: multiplying $\eta$ by a positive coefficient $c$ multiplies
$\sqrt{|\det\eta|}$ by $c^{3/2}$. Hence $\operatorname{vol}_{g_\rho}=\mu_\rho$
and its total volume is one. Positive conformal scaling preserves time
orientation and chronology.

For $n\geq0$, sample $P_1,\ldots,P_n$ independently from $\mu_\rho$ and
form the matrix
\[
 A_n(P)_{ij}=\mathbf 1_{\{P_i\ll P_j\}}.
\]
Write $\nu_n^\rho$ for its law. This is a Borel pushforward because the
inequality in~\eqref{eq:chronology} is open. Let $\bar\nu_n^\rho$ be the
pushforward under simultaneous relabeling of rows and columns. This
forgets vertex names and retains the direction of the order; it does
not identify an order with its time reverse.

For $R\in\Otwo$, let $T_R(t,z)=(t,Rz)$. The proposed discrepancy is
\begin{equation}\label{eq:distance}
 d_{\Otwo}(\rho,\sigma)=
 \inf_{R\in\Otwo}\sup_{p\in\D}|\rho(p)-\sigma(T_Rp)|.
\end{equation}
The group includes reflections. It does not include nonlinear conformal
maps of the diamond.

\subsection{Main statement}
For $|a|<1$, define
\begin{align}
 D_a(t,z)&=(1+at)^2-a^2|z|_2^2,\label{eq:denominator}\\
 N_a(t,z)&=(t+a)(1+at)-a|z|_2^2,\nonumber\\
 \Omega_a(t,z)&=\frac{1-a^2}{D_a(t,z)},\nonumber\\
 F_a(t,z)&=\left(\frac{N_a(t,z)}{D_a(t,z)},\,
                  \frac{(1-a^2)z}{D_a(t,z)}\right).\label{eq:flow}
\end{align}
The formulas are used on the positive-denominator neighborhood of the
closed diamond. Fix from now on
\begin{equation}\label{eq:density}
 a_*=\frac1{100},\qquad \rho_0=1,\qquad
 \rho_* =\Omega_{-a_*}^{\,3}.
\end{equation}

\begin{theorem}\label{thm:main}
The densities $\rho_0,\rho_*$ belong to $\K$. For every $n\geq0$,
\[
 \nu_n^{\rho_*}=\nu_n^{\rho_0},\qquad
 \bar\nu_n^{\rho_*}=\bar\nu_n^{\rho_0}.
\]
Nevertheless $d_{\Otwo}(\rho_0,\rho_*)>0$. The map $F_{-a_*}$ is a
smooth time-oriented diffeomorphism of $\D$ satisfying
\begin{equation}\label{eq:isometry}
 g_{\rho_*}=F_{-a_*}^{\,*}g_{\rho_0}.
\end{equation}
Both maps and their inverse formulas extend smoothly to neighborhoods
of the closed diamond, including the spatial axis and boundary points.
\end{theorem}

Sections~\ref{sec:flow}--\ref{sec:separation} prove the theorem. The
isometry in~\eqref{eq:isometry} identifies the two geometries exactly.
The positive discrepancy is a failure of the proposed coordinate gauge
to identify two descriptions of that geometry.

\section{The conformal diamond flow}\label{sec:flow}
\begin{lemma}\label{lem:automorphism}
For $|a|<1$, $F_a$ maps the open and closed diamond bijectively to
themselves, with inverse $F_{-a}$. It is smooth on an open neighborhood
of $\overline\D$.
\end{lemma}
\begin{proof}
Use radial null variables $u=t+r$ and $v=t-r$. Membership in
$\overline\D$ is equivalent to $|u|,|v|\leq1$, and
\[
 D_a=(1+au)(1+av)\geq(1-|a|)^2>0.
\]
Thus $U_a=\{D_a>0\}$ is an open neighborhood of $\overline\D$ and
the Cartesian rational formulas are smooth there. There is no division
by $r$ in~\eqref{eq:flow}.

The function $f_a(s)=(s+a)/(1+as)$ satisfies
\[
 f_a'(s)=\frac{1-a^2}{(1+as)^2}>0,\qquad
 f_a(\pm1)=\pm1,\qquad f_{-a}(f_a(s))=s.
\]
Directly from~\eqref{eq:flow}, positivity of $\Omega_a$ and $r\geq0$,
\begin{equation}\label{eq:null}
 t' + r'=f_a(u),\qquad t'-r'=f_a(v),\qquad r'=\Omega_a r.
\end{equation}
This proves preservation of the open and closed diamond. When $r>0$,
the spatial direction is unchanged, so applying $f_{-a}$ to both null
variables restores $t,r,z$. At $r=0$ the spatial coordinates stay zero
and the same inverse identity restores $t$. Applying the argument
with $-a$ proves both inverse compositions and surjectivity. The inverse
also has a positive-denominator neighborhood, proving the asserted
smoothness up to the boundary in the ambient-coordinate sense.
\end{proof}

\begin{lemma}\label{lem:conformal}
For $p\in\overline\D$ and $|a|<1$,
\begin{equation}\label{eq:derivative}
 (\dd F_a)_p^{\mathsf T}\eta(\dd F_a)_p=\Omega_a(p)^2\eta,
 \qquad \det(\dd F_a)_p=\Omega_a(p)^3>0.
\end{equation}
For $p,q\in\overline\D$, the separation identity is
\begin{equation}\label{eq:separation}
 \eta(F_a(q)-F_a(p),F_a(q)-F_a(p))
 =\Omega_a(p)\Omega_a(q)\eta(q-p,q-p).
\end{equation}
\end{lemma}
\begin{proof}
Away from the axis use polar coordinates, so that
$\eta=-\dd u\,\dd v+r^2\dd\theta^2$ (with the symmetric interpretation
of the mixed term). The two null derivatives have product
$f_a'(u)f_a'(v)=\Omega_a^2$, and $r'^2=\Omega_a^2r^2$ while the
angle stays fixed. This proves the bilinear identity there. Cartesian
smoothness and continuity extend it to the axis and boundary.

Taking determinants yields $(\det\dd F_a)^2=\Omega_a^6$. For a fixed
point of $\overline\D$, replace $a$ by $sa$, $0\leq s\leq1$. The
determinant is continuous and never zero, and is one at $s=0$.
Its sign is therefore positive throughout, proving the determinant
formula. This sign step is needed for the density transport below.

For completeness,~\eqref{eq:separation} is an identity of rational
Cartesian functions. Multiply by $D_a(p)^2D_a(q)^2$, insert
the three numerators in~\eqref{eq:flow}, and expand: the resulting
polynomial is
\[
 (1-a^2)^2D_a(p)D_a(q)
 \bigl(-(q_t-p_t)^2+(q_x-p_x)^2+(q_y-p_y)^2\bigr).
\]
Both denominators are positive, giving~\eqref{eq:separation}.
Appendix~\ref{app:algebra} gives a Cartesian derivative certificate
which also avoids polar coordinates altogether.
\end{proof}

\begin{proposition}\label{prop:chronology}
For $|a|<1$ and $p,q\in\overline\D$,
$p\ll q$ if and only if $F_a(p)\ll F_a(q)$.
\end{proposition}
\begin{proof}
Suppose $p\ll q$. Along the continuous deformation $s\mapsto F_{sa}$,
equation~\eqref{eq:separation} keeps the squared Lorentz separation
strictly negative. The transformed time difference cannot be zero,
because a vector with zero time component has nonnegative Lorentz square.
It starts positive at $s=0$, hence remains positive at $s=1$ by the
intermediate value theorem. Negative square and positive time difference
together give~\eqref{eq:chronology} for the images. Apply the same
argument to $F_{-a}$ and use Lemma~\ref{lem:automorphism} for the converse.
This proves orientation preservation for arbitrary pairs, including
pairs with different spatial directions.
\end{proof}

\section{Density bounds and actual measure transport}\label{sec:density}
\begin{lemma}\label{lem:bounds}
The function $\rho_*$ is smooth on a neighborhood of $\overline\D$ and
satisfies the bounds and Euclidean Lipschitz condition in~\eqref{eq:class}.
\end{lemma}
\begin{proof}
Smoothness follows from $D_{-a_*}>0$. Factoring in null variables gives
\[
 (99/100)^2\leq D_{-a_*}(p)\leq(101/100)^2,
 \qquad
 (99/101)^3\leq\rho_*(p)\leq(101/99)^3.
\]
The latter interval lies inside $[1/2,3/2]$, as follows by multiplying
the positive integer denominators.

Here is a direct global Lipschitz estimate using Euclidean distances.
Every coordinate of a point of $\overline\D$ has absolute value at most
one. Hence, for any such $p,q$ and coordinate $i$,
\[
 |p_i-q_i|\leq\|p-q\|_2,\qquad
 |p_i^2-q_i^2|\leq2\|p-q\|_2.
\]
Expansion of $D_{-a_*}=1-2a_*t+a_*^2(t^2-x^2-y^2)$ therefore gives
\begin{equation}\label{eq:denom-lip}
 |D_{-a_*}(p)-D_{-a_*}(q)|
 \leq(2a_*+6a_*^2)\|p-q\|_2
 \leq\frac3{100}\|p-q\|_2.
\end{equation}
If $9/10\leq d,e\leq11/10$ and $0\leq c\leq1$, then
\begin{align*}
 \left|\frac c{d^3}-\frac c{e^3}\right|
 &=\frac{c|d-e|(d^2+de+e^2)}{d^3e^3}\\
 &\leq\frac{363/100}{(9/10)^6}|d-e|\leq7|d-e|.
\end{align*}
Take $d=D_{-a_*}(p)$, $e=D_{-a_*}(q)$ and $c=(1-a_*^2)^3$.
Their stated bounds satisfy these hypotheses. Combining this with
\eqref{eq:denom-lip} gives a Lipschitz constant at most $21/100<2$.
The formal class theorem records the required constant two.
\end{proof}

\begin{proposition}\label{prop:transport}
For every $|a|<1$,
\begin{equation}\label{eq:transport}
 (F_a)_*\mu_0=\Omega_{-a}^{\,3}\mu_0.
\end{equation}
In particular $\rho_*$ is normalized and belongs to $\K$, as does $1$.
\end{proposition}
\begin{proof}
The change-of-variables formula for the diffeomorphism $F_a:\D\to\D$
and Lemma~\ref{lem:conformal} give, for every nonnegative Borel function $h$,
\begin{align*}
 \int_\D h(F_a(p))\,\dd\mu_0(p)
 &=\frac1V\int_\D h(q)|\det\dd F_{-a}(q)|\,\dd q\\
 &=\int_\D h(q)\Omega_{-a}(q)^3\,\dd\mu_0(q).
\end{align*}
This proves equality of the actual measures. Setting $h=1$ proves
normalization, rather than assuming it as an additional hypothesis.
Lemma~\ref{lem:bounds} supplies the remaining conditions for $a=a_*$.
The constant function one has all class properties immediately
from~\eqref{eq:volume}.
\end{proof}

\section{Finite order laws}\label{sec:laws}
\begin{proposition}\label{prop:laws}
For every $n\geq0$, both finite-order laws in Theorem~\ref{thm:main} agree.
\end{proposition}
\begin{proof}
Let $P_1,\ldots,P_n$ be independent with law $\mu_0$ and define
$Q_i=F_{a_*}(P_i)$. Deterministic coordinatewise transport preserves
independence, and Proposition~\ref{prop:transport} makes each $Q_i$
have law $\mu_{\rho_*}$. Equivalently,
\[
 (F_{a_*}^{\times n})_*(\mu_0^{\otimes n})
 =\mu_{\rho_*}^{\otimes n}.
\]
Proposition~\ref{prop:chronology} gives $A_n(Q)=A_n(P)$ pointwise for
every sample in $\D^n$. Taking laws proves the labeled equality.
Pushing both equal measures through the same relabeling quotient gives
the unlabeled equality. At $n=0$ both laws are the unit mass on the
empty order. No estimation, asymptotic limit or time-reversal quotient
is used.
\end{proof}

\section{Coordinate separation and geometric isometry}\label{sec:separation}
\begin{proposition}\label{prop:positive}
The discrepancy in~\eqref{eq:distance} is positive. More explicitly,
\begin{equation}\label{eq:anchor}
 d_{\Otwo}(1,\rho_*)\geq
 \left(\frac{39996}{39601}\right)^3-1>0.
\end{equation}
\end{proposition}
\begin{proof}
The density depends on $z$ only through $|z|_2^2$, so
$\rho_*\circ T_R=\rho_*$ for every $R\in\Otwo$. Each supremum in
\eqref{eq:distance} is thus the same number. It is finite because
$|\rho_*-1|\leq1/2$. At the interior point $p_*=(1/2,0,0)$,
\[
 \rho_*(p_*)=
 \left(\frac{1-1/10000}{(1-1/200)^2}\right)^3
 =\left(\frac{39996}{39601}\right)^3>1.
\]
The supremum bounds this value minus one for every $R$. Taking the
infimum proves~\eqref{eq:anchor}. This uses an interior point and hence
does not depend on whether boundary values enter the coordinate norm.
\end{proof}

\begin{proposition}\label{prop:isometry}
The map $H=F_{-a_*}$ is a smooth time-oriented isometry from
$(\D,g_{\rho_*})$ onto $(\D,g_1)$.
\end{proposition}
\begin{proof}
The previous results give its smooth inverse and time orientation.
Since $\Omega_{-a_*}>0$,
\[
 \rho_*^{2/3}=(\Omega_{-a_*}^3)^{2/3}
 =\Omega_{-a_*}^2.
\]
Using the derivative pullback identity in~\eqref{eq:derivative},
\[
 H^*g_1=V^{-2/3}H^*\eta
       =V^{-2/3}\Omega_{-a_*}^2\eta=g_{\rho_*}.
\]
This equality holds for every tangent-vector pair at every point,
including the ambient extension to $\overline\D$. It proves the claimed
isometry directly. The sampling measures transform in the compatible
direction $H_*\mu_{\rho_*}=\mu_0$.
\end{proof}

These propositions complete Theorem~\ref{thm:main}. In particular the
nonconstant coordinate coefficient in $g_{\rho_*}$ still describes a
flat geometry. A coordinate formula alone does not distinguish
Lorentzian isometry classes.

Define, for either the labeled or unlabeled directed laws,
\[
 \Delta_N(\rho,\sigma)=\max_{1\leq n\leq N}
          \TV(\nu_n^\rho,\nu_n^\sigma).
\]
\begin{corollary}\label{cor:inverse}
There is no function $b:\mathbb N\times[0,1]\to[0,\infty)$ with
$b(N,0)\to0$ such that, for every pair in $\K$ and all $N\geq1$,
\[
 d_{\Otwo}(\rho,\sigma)\leq b(N,\Delta_N(\rho,\sigma)).
\]
The assertion holds for either choice of finite-order laws.
\end{corollary}
\begin{proof}
Insert $\rho=1$ and $\sigma=\rho_*$. Proposition~\ref{prop:laws}
gives $\Delta_N=0$ for every $N$, while Proposition~\ref{prop:positive}
gives a fixed positive lower bound for the left side. Letting $N$
tend to infinity is a contradiction.
\end{proof}
This corollary is an elementary prose consequence of the certified
example. It asserts failure of this particular coordinate discrepancy;
it supplies no inverse rate for a corrected gauge.

\section{Comparison with established work}\label{sec:prior}
\subsection{The flow is established}
Casini, Huerta and Myers give the causal-diamond flow in
equation (2.13) of their arXiv v2 PDF~\cite{CHM}; their discussion
cites Hislop and Longo~\cite{HL} for earlier double-cone modular
geometry. To identify their formula with ours, set the ball radius
to one and $a=\tanh(\pi s)$. Their $x^+=r+t$ becomes $u$,
and their $x^-=r-t$ becomes $-v$.
For $k=e^{-2\pi s}$, elementary algebra gives
\[
 \frac{(1+w)-k(1-w)}{(1+w)+k(1-w)}
 =\frac{w+a}{1+aw},\qquad a=\frac{1-k}{1+k}.
\]
The opposite sign in the $x^-$ exponential gives the same $f_a$
after replacing $x^-$ by $-v$. Equations~\eqref{eq:null} then give
exactly~\eqref{eq:flow}. Thus this manuscript does not introduce a
new conformal transformation. Its parameter $a_*$ selects a small
member of that established family.

\subsection{Probability-law equivalence}
Janson's Theorem 7.1~\cite{Janson} relates equivalent poset kernels,
their finite random-poset laws and measure-preserving representations.
Our deterministic kernel is $W(p,q)=\mathbf1_{p\ll q}$.
Propositions~\ref{prop:transport} and~\ref{prop:chronology} supply an
explicit measure-preserving order isomorphism between the two measured
spaces. The coupling proof is the direct invariance direction of this
general framework; it does not provide a new abstract equivalence theorem.

\subsection{Reconstruction up to isometry}
Braun's Theorem 1.4~\cite{Braun} reconstructs finite-volume, causally
continuous, future-chronocomplete spacetimes up to smooth isometry from
all finite labeled chronological adjacency laws. The weighted version
and Remark 1.6 explain the corresponding conformal rescaling.
Our explicit pullback identity gives the isometry directly. Positive
coordinate distance is therefore compatible with that reconstruction
target. We use no unverified application of its converse or its
geometric hypotheses. The equality of unlabeled laws here follows
from the explicit labeled coupling and makes no claim to solve a
general reconstruction conjecture for unlabeled data.

\subsection{Contribution and limits of the comparison}
The flow, order-law invariance under transport, and geometric
interpretation are established ingredients. This paper packages their
explicit small-parameter realization inside~\eqref{eq:class}, proves
the coordinate discrepancy and pullback together, and exposes the
formal proof and acceptance evidence. It can serve as a checked test
case for proposed inverse distances and gauge normalizations.
The comparison covers the displayed flow and the cited reconstruction
and kernel-equivalence results. It is not an exhaustive priority search
for this exact bounded-density example or its formalization; no first-result
claim is made.

\section{Formal statements and reproducibility}\label{sec:formal}
The proof is implemented in Lean 4.30.0 with mathlib revision
\begin{center}
\code{c5ea00351c28e24afc9f0f84379aa41082b1188f}.
\end{center}
Cartesian points use
\code{EuclideanSpace} over three real coordinates. The norm is genuinely
Euclidean, chronology uses the spatial radius, and all sampling laws
are actual measure-theoretic products and pushforwards.

The main export is
\begin{center}
\code{QuantyraNullCone.higher\_dimensional\_gauge\_counterexample3}.
\end{center}
It has no input hypotheses: the explicit parameter is part of the
definitions. Its conjunction gives both density-class memberships,
all labeled and unlabeled law equalities, a positive infimum over
all spatial linear isometries, and the metric pullback for every
point and tangent-vector pair. Smooth automorphism and chronology
are also exported separately. The following map identifies the key
statements, all in namespace \code{QuantyraNullCone}:
\begin{itemize}
\raggedright
\item Volume and probability: \code{lorentz\_diamond\_volume3},
\code{flat\_density\_class3}, in \code{LorentzVolume.lean} and
\code{LorentzFlat.lean}.
\item Lemma~\ref{lem:automorphism}: \code{lorentz\_flow\_smooth\_automorphism3},
in \code{LorentzSmooth.lean}, with inverse and preservation dependencies.
\item Lemma~\ref{lem:conformal}: \code{hasFDerivAt\_lorentz\_flow3},
\code{lorentz\_flow\_derivative\_det3},
\code{lorentz\_flow\_derivative\_bilinear3} and
\code{lorentz\_flow\_separation3}, in \code{LorentzDerivative.lean},
\code{LorentzJacobian.lean} and \code{LorentzChronology.lean}.
\item Proposition~\ref{prop:chronology}: \code{lorentz\_flow\_chronological\_iff3},
in \code{LorentzChronology.lean}.
\item Lemma~\ref{lem:bounds} and Proposition~\ref{prop:transport} at $a_*$:
\code{gauge\_density\_bounds3}, \code{gauge\_density\_lipschitz3},
\code{gauge\_density\_class3}, \code{gauge\_forward\_transport3},
in the \code{LorentzGauge} modules and \code{LorentzTransport.lean}.
\item Proposition~\ref{prop:laws}: \code{gauge\_order\_laws\_equal3} and
\code{gauge\_unlabeled\_order\_laws\_equal3}, in \code{LorentzOrderLaws.lean}.
\item Proposition~\ref{prop:positive}: \code{gauge\_anchor\_density\_gt3} and
\code{gauge\_o2\_distance\_positive3}, in \code{LorentzSpatialDistance.lean}.
\item Proposition~\ref{prop:isometry}: \code{gauge\_metric\_isometry3},
in \code{LorentzIsometry.lean}.
\end{itemize}

The general-$a$ transport statement, displayed sharper density and
Lipschitz constants, explicit rational anchor lower bound,
Corollary~\ref{cor:inverse}, and literature parameter conversion are
proved in the exposition. The accepted Lean exports give the fixed
$a_*$ transport, class bounds, Lipschitz constant two and strict positive
distance; we do not claim separate exports for every displayed refinement.
Likewise, the interpretation of the coordinate bilinear pullback as a
manifold isometry and the determinant calculation of metric volume are
spelled out here; the Lean theorem uses the actual Cartesian derivative
and bilinear metric rather than a packaged Lorentzian-manifold structure.

The accepted proof/evidence revision is
\href{https://github.com/Quantyra/quantyra-null-cone-lean/tree/74c1f743d085c63b45ac2bf30008ad5d29b98fbc}{\code{74c1f743d085c63b45ac2bf30008ad5d29b98fbc}}.
Authoritative verification ran on GCP under
\code{space-lorentz-acceptance-20261008T071627Z-221cd8}.
The root build completed 3001 jobs and 295 selected type/axiom audits
across the full repository, including 76 exports added by this campaign.
This count includes other research modules. No owned warnings or added
axioms were reported. The selected dependencies use only
\code{propext}, \code{Classical.choice} and \code{Quot.sound}.
Source hashes and all nine pinned dependency identities were checked
before and after execution. The raw logs, receipt and failed-development
ledger are retained in the repository's \code{evidence/gcp} directory.

Every development and acceptance Lean invocation for this campaign ran
on GCP. Hosted CI is supplementary. This manuscript reuses that acceptance
after checking the relevant source identities; it requires no new Lean
compilation. Its symbolic and rational algebra checks supplement the
proof audit and do not certify analytic or probabilistic claims by testing.
The manuscript directory contains build instructions, a formal statement
map, primary-source comparison and a review receipt.

\section{Implications for a corrected inverse problem}
A time-oriented conformal diffeomorphism $H:\D\to\D$ with
$H^*\eta=\Omega_H^2\eta$ acts on volume-normalized density descriptions
by
\begin{equation}\label{eq:gauge-action}
 \rho_H(p)=|\det\dd H_p|\,\sigma(H(p))
          =\Omega_H(p)^3\sigma(H(p)).
\end{equation}
Whenever the densities in question meet the required regularity,
$g_{\rho_H}=H^*g_\sigma$. Omitting the Jacobian factor would transform
scalar values, rather than the sampling-volume density. In the example,
$H=F_{-a_*}$ and $\sigma=1$.

An inverse problem for coefficients must therefore justify a normalization
that removes this freedom, or use an error quantity on the appropriate
geometric equivalence classes. The full conformal action need not preserve
the fixed bounds in $\K$. Moreover, nonconstant Jacobian weights mean
that simply taking an infimum of coefficient differences under that
action does not automatically yield a symmetric metric. A corrected
definition requires its own invariance and stability analysis.

This example does not supply a finite reconstruction lemma for that
corrected problem. The quantitative higher-dimensional inverse question
remains open in this project. It also does not alter the separate
$1+1$ theorem, whose prescribed null-coordinate marginal normalization
is a different model assumption. Both metrics here represent the same
geometry, so there is no conclusion that two distinct physical spacetimes
are indistinguishable by the observations.

The work was developed with OpenAI Codex assistance under the author's
research direction. Quantyra's review workflow is open-source,
decentralized and informal: readers can inspect, use, test and improve
the public artifacts. Specialist independent review is not a publication
prerequisite; no external review or adoption is claimed by this manuscript.

\appendix
\section{Cartesian algebra certificate}\label{app:algebra}
Let $c=1-a^2$, $D=D_a(p)$, $N=N_a(p)$ and let the numerator vector be
$A=(N,cx,cy)$. Its derivatives are
\[
 \nabla D=(2a(1+at),-2a^2x,-2a^2y),\qquad
 \nabla N=(1+a^2+2at,-2ax,-2ay).
\]
The quotient rule gives
\[
 (\dd F_a)_{ij}=\frac{B_{ij}}{D^2},\qquad
 B=D\,\dd A-A(\nabla D)^{\mathsf T}.
\]
With $E=\operatorname{diag}(-1,1,1)$, expansion gives the polynomial
identities
\[
 B^{\mathsf T}EB=c^2D^2E,\qquad \det B=c^3D^3.
\]
Division by $D^4$ and $D^6$ respectively recovers both identities
in~\eqref{eq:derivative}. These are exactly Cartesian identities,
including $x=y=0$; no limiting polar computation is needed.
The supplementary symbolic script checks these polynomials, the
inverse and separation numerators, the null-coordinate formulas and
the literature parameter conversion. Exact rational checks cover the
class constants and the interior anchor. Analytic domain restrictions,
smoothness and probability transport are justified in the proofs above
and by the mapped Lean statements, rather than by symbolic cancellation.

\begin{thebibliography}{9}
\bibitem{CHM} H. Casini, M. Huerta and R. C. Myers,
\emph{Towards a derivation of holographic entanglement entropy},
JHEP \textbf{05} (2011), 036.
\href{https://arxiv.org/abs/1102.0440v2}{arXiv:1102.0440v2}.
\bibitem{HL} P. D. Hislop and R. Longo,
\emph{Modular structure of the local algebras associated with the free
massless scalar field theory}, Commun. Math. Phys. \textbf{84} (1982),
71--85. \href{https://doi.org/10.1007/BF01208372}{doi:10.1007/BF01208372}.
\bibitem{Janson} S. Janson,
\emph{Poset limits and exchangeable random posets},
\href{https://arxiv.org/abs/0902.0306v1}{arXiv:0902.0306v1} (2009),
Theorem 7.1.
\bibitem{Braun} M. Braun,
\emph{Spacetime reconstruction by order and number},
\href{https://arxiv.org/abs/2507.01907v1}{arXiv:2507.01907v1} (2025).
\bibitem{Fredriksen} D. E. Fredriksen,
\emph{Quantitative reconstruction from finite causal orders in a conformal
diamond}, version 0.4.0 (2026), Quantyra Inc.
\href{https://doi.org/10.5281/zenodo.23247720}{doi:10.5281/zenodo.23247720}.
\end{thebibliography}
\end{document}
